\documentclass[11pt,oneside]{book}
\usepackage[T1]{fontenc}
\usepackage[utf8]{inputenc}
\usepackage{lmodern}
\usepackage[a4paper,margin=25mm]{geometry}
\usepackage{amsmath,amssymb,enumitem,microtype,needspace}
\usepackage[hidelinks]{hyperref}
\hypersetup{pdftitle={Mathematical Foundations of Statistics and Learning: Book Plan}}
\pagestyle{plain}
\setlength{\parindent}{0pt}
\setlength{\parskip}{4pt}
\setcounter{tocdepth}{1}
\setlist[itemize]{leftmargin=1.5em,itemsep=1pt,topsep=3pt}
\newcounter{planchapter}
\newcommand{\PlanPart}[2]{\Needspace{18\baselineskip}\section*{Part #1. #2}\addcontentsline{toc}{section}{Part #1. #2}}
\newcommand{\PlanChapter}[2]{\Needspace{14\baselineskip}\refstepcounter{planchapter}\subsection*{\theplanchapter. #2}\label{plan:#1}}
\newenvironment{Outline}{\begin{enumerate}[label=\theplanchapter.\arabic*,leftmargin=2.8em,itemsep=1pt,topsep=2pt]}{\end{enumerate}}
\newcommand{\Target}[1]{\textit{Principal target.} #1\par}
\newcommand{\Examples}[1]{\textit{Examples.} #1\par}
\title{Mathematical Foundations of Statistics and Learning\\[6pt]\large Book plan: the route to transformers and language models}
\author{Working plan}
\date{11 November 2026}

\begin{document}
\hypersetup{pageanchor=false}
\maketitle
\clearpage
\hypersetup{pageanchor=true}
\frontmatter
\tableofcontents

\chapter*{Scope and conventions}
\addcontentsline{toc}{chapter}{Scope and conventions}

\textbf{First overall goal.} Develop a mathematical account of autoregressive
transformers and language models: statistical interpretation, architecture,
training objectives, differentiation, inference algorithms, and conditional
theoretical guarantees. The proposal contains eight parts and 34 chapters.
Titles and numbering remain provisional.

\textbf{Agreed background.} Standard measure-theoretic probability, including
its core measure theory, may be assumed. State imported results with exact
hypotheses when used. Specialized statistics, information theory, optimization,
and learning theory are developed in the text. Record borderline imports
explicitly.

\textbf{Additional background proposed for discussion.} Elementary
finite-dimensional linear algebra, real analysis, and multivariable calculus.
This boundary is not yet agreed. Develop specialized tools---matrix
differentiation, convex duality, reproducing kernels, approximation results---where
needed. Introduce definitions at first use; later chapters develop their theory
without redefining them.

\textbf{Form.} Definitions, results, concise proofs, numerous worked examples,
counterexamples, and exercises; minimal discussion. Plan a section or subsection,
then draft definitions, statements, and proofs in separate passes. The author's
study drives revision. Python is deferred to a later phase.

\textbf{Result status.} The targets below specify intended coverage, not
theorems established in this project. Each requires precise hypotheses during
unit planning. Architectural identities, statistical guarantees, empirical
observations, and open questions remain distinct. In particular, practical
neural training is not assumed to attain global optima.

\textbf{Breadth and order.} Classical inference, kernels, and latent-variable
methods retain a place even where they are not direct transformer prerequisites.
The dependency graph records what each result needs. The project may eventually
span several volumes; the eight-part structure is a mathematical organization,
not a promise of a short physical book.

\mainmatter
\chapter*{Proposed contents}
\addcontentsline{toc}{chapter}{Proposed contents}

\PlanPart{I}{Statistical models and decisions}

\PlanChapter{models}{Statistical models and observations}
\begin{Outline}
\item Observation spaces, models, and parametrized families.
\item Sampling schemes and product experiments.
\item Parameters, statistical functionals, and identifiability.
\item Domination, densities, and likelihoods.
\end{Outline}
\Target{Characterize identifiable targets through their underlying laws;
derive likelihoods for specified observation schemes.}
\Examples{Bernoulli, Gaussian, and uniform endpoint families; redundant
parametrizations; independent versus dependent observations.}

\PlanChapter{decisions}{Loss, risk, and decision rules}
\begin{Outline}
\item Actions and deterministic or randomized rules.
\item Loss functions, risk, and comparison of procedures.
\item Bayes risk and optimal actions in elementary models.
\item Minimax risk, admissibility, and existence issues.
\end{Outline}
\Target{Derive optimal decisions in finite models and under square loss;
distinguish optimal values from attained optimal rules.}
\Examples{Mean estimation; asymmetric classification costs; ties and nonattainment.}

\PlanChapter{information}{Information, divergence, and scoring rules}
\begin{Outline}
\item Entropy, cross-entropy, and logarithmic loss.
\item Relative entropy and divergence inequalities.
\item Conditional information, chain rules, and data processing.
\item Proper scoring rules and probability estimation.
\end{Outline}
\Target{Prove information inequalities and the cross-entropy decomposition;
characterize minimizers of logarithmic and quadratic probability scores.}
\Examples{Finite laws; Gaussian divergence; incompatible supports and infinite values.}

\PlanChapter{sufficiency}{Statistics, sufficiency, and information reduction}
\begin{Outline}
\item Statistics and induced experiments.
\item Sufficiency and conditional formulations.
\item Factorization criteria under explicit assumptions.
\item Minimal sufficiency, completeness, and ancillarity.
\end{Outline}
\Target{Establish sufficient reductions in dominated models; distinguish
sufficiency from identifiability and completeness.}
\Examples{Binomial counts; Gaussian summaries; order statistics; information discarded by a statistic.}

\PlanPart{II}{Statistical inference}

\PlanChapter{estimation}{Point estimation and finite-sample optimality}
\begin{Outline}
\item Estimators, bias, variance, and mean squared error.
\item Maximum likelihood and other estimating constructions.
\item Rao--Blackwellization and complete sufficient statistics.
\item Fisher information and information lower bounds.
\end{Outline}
\Target{Prove Rao--Blackwell and Lehmann--Scheff\'e results and a regular
Cram\'er--Rao bound; expose their required hypotheses.}
\Examples{Bernoulli and normal estimation; biased improvements; boundary
maxima; parameter-dependent support.}

\PlanChapter{testing}{Tests and confidence procedures}
\begin{Outline}
\item Hypotheses, randomized tests, size, and power.
\item Neyman--Pearson theory and composite hypotheses.
\item Confidence sets and inversion of tests.
\item Multiple decisions and error criteria.
\end{Outline}
\Target{Prove the Neyman--Pearson lemma and testing--confidence duality;
derive exact procedures in tractable models.}
\Examples{Binomial and Gaussian tests; discrete coverage; absence of a uniformly most powerful test.}

\PlanChapter{bayes}{Bayesian inference}
\begin{Outline}
\item Priors, posterior laws, and existence conditions.
\item Conjugate models and posterior prediction.
\item Posterior decisions and credible sets.
\item Prior sensitivity and improper-prior pitfalls.
\end{Outline}
\Target{Derive posterior and predictive laws and Bayes decisions;
distinguish posterior probability from repeated-sampling coverage.}
\Examples{Beta--Bernoulli; normal conjugacy; alternative priors; an improper posterior.}

\PlanChapter{asymptotics}{Consistency and asymptotic inference}
\begin{Outline}
\item Consistency, continuous transformations, and the delta method.
\item Compact-parametric uniform laws from probability prerequisites.
\item Consistency and asymptotic normality of regular estimators.
\item Likelihood-based asymptotic procedures and their failures.
\end{Outline}
\Target{Prove argmin consistency and regular asymptotic normality. Establish
needed uniform laws here without relying on later chapters.}
\Examples{Sample moments; normal likelihood; boundary models; pointwise
convergence insufficient for an argmin argument.}

\PlanChapter{lower-bounds}{Minimax inference and lower bounds}
\begin{Outline}
\item Minimax risk, rates, and finite-sample comparisons.
\item Reduction from estimation to testing.
\item Two-point and multiple-hypothesis lower bounds.
\item Matching upper and lower bounds.
\end{Outline}
\Target{Derive Le Cam and Fano methods in stated settings and establish
optimal rates in elementary models.}
\Examples{Bernoulli mean; Gaussian location; finite packing constructions.}

\PlanPart{III}{Computation, differentiation, and optimization}

\PlanChapter{computation}{Algorithms, cost models, and numerical stability}
\begin{Outline}
\item Algorithms, representations, and input size.
\item Arithmetic cost, bit cost, memory, and asymptotic notation.
\item Conditioning and floating-point error.
\item Stable elementary statistical operations.
\end{Outline}
\Target{State computational bounds relative to an explicit model;
separate exact identities from numerical approximations.}
\Examples{Matrix multiplication; ill-conditioned least squares; stable log-sum-exp and softmax.}

\PlanChapter{differentiation}{Differentiation and computational graphs}
\begin{Outline}
\item Derivatives as linear maps; matrix-valued functions.
\item Chain rules, Jacobians, and adjoints.
\item Forward and reverse automatic differentiation.
\item Nondifferentiability and derivative conventions.
\end{Outline}
\Target{Prove reverse accumulation correct on finite computation graphs;
state time and memory costs under explicit assumptions.}
\Examples{Least squares; log-likelihood gradients; shared parameters; piecewise-linear activations.}

\PlanChapter{convexity}{Convex analysis and duality}
\begin{Outline}
\item Convex sets, functions, and supporting inequalities.
\item Smoothness, strong convexity, and subgradients.
\item Conjugacy, constrained problems, and duality.
\item Optimality conditions and proximal maps.
\end{Outline}
\Target{Prove finite-dimensional optimality and duality results with
constraint qualifications; derive basic proximal operators.}
\Examples{Quadratics; entropy; norm penalties; a duality gap.}

\PlanChapter{deterministic-optimization}{Deterministic optimization}
\begin{Outline}
\item Descent lemmas and gradient methods.
\item Convex and strongly convex rates.
\item Projected, proximal, and second-order methods.
\item Nonconvex stationarity and limitations.
\end{Outline}
\Target{Prove representative convergence rates; distinguish stationarity
from global optimization.}
\Examples{Quadratic and logistic objectives; constrained least squares; saddles and flat directions.}

\PlanChapter{stochastic-optimization}{Stochastic optimization}
\begin{Outline}
\item Stochastic gradients and sampling assumptions.
\item Step sizes, minibatches, and convex guarantees.
\item Smooth nonconvex stationarity bounds.
\item Momentum and adaptive methods: qualified results.
\end{Outline}
\Target{Prove SGD guarantees under explicit noise and regularity assumptions.
Give adaptive methods only the guarantees actually established.}
\Examples{Noisy quadratics; finite-sum losses; biased gradients; unsuitable step sizes.}

\PlanPart{IV}{Statistical learning theory}

\PlanChapter{learning}{Prediction and empirical risk minimization}
\begin{Outline}
\item Prediction problems, hypothesis classes, and population risk.
\item Bayes predictors for regression and classification.
\item Empirical risk minimization and excess risk.
\item Approximation, estimation, and optimization error.
\end{Outline}
\Target{Derive risk decompositions and elementary empirical-risk solutions;
distinguish optimization from statistical performance.}
\Examples{Constant predictors; finite classes; thresholds; interpolation and unattained minima.}

\PlanChapter{concentration}{Concentration and uniform convergence}
\begin{Outline}
\item Exponential-moment methods.
\item Hoeffding, Bernstein, and bounded-difference bounds.
\item Finite classes, coverings, and uniform deviations.
\item Generalization of empirical minimizers.
\end{Outline}
\Target{Prove concentration and a finite-class excess-risk bound; record
any overlap with the chosen probability prerequisites.}
\Examples{Empirical means; rare events; bounded versus heavy-tailed losses.}

\PlanChapter{capacity}{Capacity, complexity, and learnability}
\begin{Outline}
\item PAC formulations and sample complexity.
\item Growth functions, VC dimension, and combinatorial bounds.
\item Rademacher complexity and symmetrization.
\item Learnability theorems and lower bounds.
\end{Outline}
\Target{Establish VC and Rademacher guarantees in precise settings, including
measurability and function-class restrictions.}
\Examples{Intervals; halfspaces; finite-dimensional predictors; an unrestricted unlearnable class.}

\PlanChapter{regularization}{Regularization, stability, and model selection}
\begin{Outline}
\item Penalized and constrained learning.
\item Stability and expected generalization.
\item Validation, selection, and repeated data reuse.
\item Surrogate losses and calibration.
\end{Outline}
\Target{Prove a stability guarantee for a regularized learner and a model-selection
bound under stated sampling assumptions.}
\Examples{Ridge; logistic classification; finite model selection; validation overfitting.}

\PlanPart{V}{Classical models and representations}

\PlanChapter{linear-models}{Linear models and regularized regression}
\begin{Outline}
\item Least squares, projections, and nonunique solutions.
\item Gaussian linear inference.
\item Ridge, shrinkage, and bias--variance calculations.
\item Sparse penalties and optimality conditions.
\end{Outline}
\Target{Derive solutions and finite-sample risk formulas under explicit
design and noise assumptions.}
\Examples{Rank deficiency; correlated features; minimum-norm interpolation; orthogonal-design thresholding.}

\PlanChapter{exponential-families}{Exponential families and generalized linear models}
\begin{Outline}
\item Natural parameters, sufficient statistics, and log partition functions.
\item Moments, curvature, and information.
\item Generalized linear models and likelihood fitting.
\item Existence, identifiability, and separation.
\end{Outline}
\Target{Prove log-partition derivative identities under valid conditions;
derive Gaussian, Bernoulli, and categorical regression.}
\Examples{Logistic, softmax, and Poisson regression; complete separation.}

\PlanChapter{kernels}{Kernels and nonparametric prediction}
\begin{Outline}
\item Positive-definite kernels and feature maps.
\item Reproducing kernel Hilbert spaces.
\item Representer theorems and kernel ridge regression.
\item Local smoothing and elementary nonparametric rates.
\end{Outline}
\Target{Construct the needed kernel spaces and prove a representer theorem;
derive bounds in a specified smoothness class.}
\Examples{Linear and Gaussian kernels; nearest neighbors; kernel smoothers.}

\PlanChapter{latent-models}{Dimension reduction and latent-variable models}
\begin{Outline}
\item PCA, low-rank approximation, and reconstruction.
\item Clustering and quantization.
\item Mixtures, label symmetry, and likelihood pathologies.
\item Variational identities and expectation--maximization.
\end{Outline}
\Target{Prove PCA optimality and EM monotonicity under stated conditions;
distinguish monotonic improvement from global optimality.}
\Examples{Gaussian mixtures; local minima; singular covariance and unbounded likelihood.}

\PlanPart{VI}{Neural networks}

\PlanChapter{neural-maps}{Neural maps and architectural structure}
\begin{Outline}
\item Tensors, affine maps, activations, and composition.
\item Multilayer, residual, and shared-parameter architectures.
\item Convolution, equivariance, and elementary symmetry.
\item Normalization and parameter symmetries.
\end{Outline}
\Target{Define architectures as maps with explicit dimensions; prove
elementary symmetry and composition properties.}
\Examples{A small multilayer network; a residual block; convolution; stabilized normalization.}

\PlanChapter{approximation}{Expressivity and approximation}
\begin{Outline}
\item Function classes, approximation metrics, and domains.
\item Piecewise-linear representations.
\item A constructive universal-approximation theorem.
\item Width, depth, and limitations.
\end{Outline}
\Target{Prove approximation for a specified network class and domain;
keep representation separate from learnability.}
\Examples{Piecewise-linear functions; finite interpolation; insufficient width.}

\PlanChapter{neural-learning}{Learning and optimizing neural models}
\begin{Outline}
\item Backpropagation as reverse differentiation.
\item Initialization, normalization, and gradient propagation.
\item Optimization in tractable special models.
\item Generalization, implicit bias, and limitations.
\end{Outline}
\Target{Derive neural gradients and selected optimization or implicit-bias
results in explicit settings, such as deep linear models.}
\Examples{Linear versus nonlinear networks; interpolation; redundancy;
valid but uninformative capacity bounds.}

\PlanPart{VII}{Sequences and transformers}

\PlanChapter{sequences}{Sequence distributions and language modeling}
\begin{Outline}
\item Alphabets, strings, prefixes, and termination conventions.
\item Conditional sequence laws and autoregressive factorization.
\item Markov, recurrent, and neural conditional models.
\item Sequence log loss, token weighting, and perplexity.
\end{Outline}
\Target{Prove normalization for specified finite-horizon or terminating models;
derive sequence likelihood and information decompositions.}
\Examples{Bigrams; a tiny vocabulary; faulty variable-length normalization;
dependent tokens within independent documents.}

\PlanChapter{tokens}{Tokenization, embeddings, and position}
\begin{Outline}
\item Encodings, tokenizers, and recoverability.
\item Pair-merging tokenization algorithms and deterministic conventions.
\item Embedding maps, output projections, and tied parameters.
\item Absolute, relative, and rotary position.
\end{Outline}
\Target{Establish encoding properties and positional identities; specify
how tokenization changes the modeled sample space.}
\Examples{Byte encodings; ambiguous segmentations; a merge table; rotary inner products.}

\PlanChapter{attention}{Attention as a mathematical operator}
\begin{Outline}
\item Queries, keys, values, and softmax attention.
\item Masks and multihead composition.
\item Equivariance, derivatives, and qualified stability results.
\item Arithmetic and memory costs.
\end{Outline}
\Target{Prove weight normalization, appropriate equivariance properties,
operator derivatives, and explicit cost bounds.}
\Examples{Uniform and concentrated attention; forbidden positions;
an empty mask row; symmetry broken by positions.}

\PlanChapter{transformers}{Transformer architectures}
\begin{Outline}
\item Encoder, decoder, and encoder--decoder constructions.
\item Residual blocks, normalization, and feedforward sublayers.
\item Causal masking and prefix dependence.
\item Parameter counts and cached evaluation.
\end{Outline}
\Target{Prove prefix causality and full-prefix/incremental equivalence
for the specified deterministic architecture in exact arithmetic.}
\Examples{A one-block decoder; pre- versus post-normalization;
cache shapes; leakage from an incorrect mask.}

\PlanPart{VIII}{Language-model training, inference, and theory}

\PlanChapter{lm-training}{Language-model estimation and training objectives}
\begin{Outline}
\item Population and empirical sequence likelihood.
\item Teacher forcing, masking, and document boundaries.
\item Gradient estimation, weighting, and optimization procedures.
\item Evaluation protocols and distribution shift.
\end{Outline}
\Target{Derive the teacher-forced likelihood and gradients; distinguish
token-weighted and document-weighted risks.}
\Examples{Padded batches; packed documents; unequal lengths; contaminated evaluation data.}

\PlanChapter{decoding}{Decoding and inference algorithms}
\begin{Outline}
\item Ancestral sampling and exact target distributions.
\item Temperature, top-$k$, and cumulative-probability truncation.
\item Greedy and beam procedures versus sequence maximization.
\item Speculative sampling and inference cost.
\end{Outline}
\Target{Prove ancestral and specified speculative sampling correct;
identify procedures that alter the target distribution.}
\Examples{A probability tree; greedy and beam failures; acceptance and residual laws.}

\PlanChapter{adaptation}{Adaptation and preference objectives}
\begin{Outline}
\item Conditional fine-tuning and distributional objectives.
\item Low-rank parameter updates.
\item Reward objectives with relative-entropy regularization.
\item Preference models and direct preference optimization.
\end{Outline}
\Target{Derive the regularized optimum and preference objective under explicit
support, normalization, and modeling assumptions.}
\Examples{Finite response sets; low-rank updates; misspecified preferences;
a failed support assumption.}
General sequential reinforcement learning and PPO belong to the later agent
extension; this chapter treats the stated single-decision objectives.

\PlanChapter{transformer-theory}{Transformer expressivity and theoretical limits}
\begin{Outline}
\item Finite-domain constructions and specified approximation results.
\item Context length, precision, and architectural restrictions.
\item In-context inference in tractable models.
\item Generalization, empirical capabilities, and open questions.
\end{Outline}
\Target{Prove selected results for precisely restricted models without
transferring them to arbitrary trained LLMs.}
\Examples{Copying and retrieval; finite lookup; tractable in-context inference;
a fixed context's limitations.}

\PlanChapter{capstone}{Capstone: a fully specified autoregressive transformer}
\begin{Outline}
\item Vocabulary, context, parameters, and the complete decoder map.
\item The probability law and likelihood objective.
\item Differentiation, cached evaluation, and sampling.
\item Computational costs and the scope of statistical guarantees.
\end{Outline}
\Target{Assemble a complete mathematical specification whose definitions
and proved properties refer to the preceding development.}
\Examples{A symbolic small-vocabulary model; end-to-end shapes, likelihood,
gradients, cache, and sampling calculations.}

\chapter*{The first unit and the dependency spine}
\addcontentsline{toc}{chapter}{The first unit and the dependency spine}

\section*{Opening unit: Section 1.1}
Plan only the opening unit to subsection depth for now:
\begin{enumerate}[label=1.1.\arabic*,leftmargin=3.4em]
\item Observation spaces and data values.
\item Statistical models as collections of probability laws.
\item Indexed families and parametrizations.
\item Worked models, counterexamples, and exercises.
\end{enumerate}
\textit{Definitions pass:} observation space, model, indexed family, parameter
space, and reparametrization. Distinguish a law from its densities and a model
from its parameter labels.

\textit{Statements pass:} invariance of a model under bijective
reparametrization; elementary examples of different parametrizations of the
same collection of laws. No general inference theorem is needed in this unit.

\textit{Proofs and examples pass:} prove the elementary facts; work through
Bernoulli, categorical, Gaussian-location, and uniform-endpoint families.
Include a redundant parametrization. Exercises ask the reader to specify
observation spaces and models, compare parametrizations, and distinguish
quantities attached to a law from quantities attached to parameter labels.
Formal identifiability follows in Section 1.3.

\section*{Logical dependencies}
These edges describe the main route, not a requirement to finish every preceding
chapter before any later chapter.
\begin{itemize}
\item Models and decisions (1--2), with information (3), support estimation
and prediction (5, 15), then learning guarantees (16--18).
\item Computation and differentiation (10--11), with convex analysis (12),
support optimization (13--14), then neural training (25, 30).
\item Linear and categorical predictors (19--20) connect statistical losses
to neural maps (23) and sequence conditionals (26).
\item Neural maps and differentiation (23, 11), with sequence and position
conventions (26--27), support attention (28) and transformers (29).
\item Information (3), sequences (26), and transformers (29) support
language-model estimation (30), decoding (31), and adaptation (32).
\item The capstone (34) assembles these constructions and cites only the
learning and optimization guarantees whose assumptions actually apply.
\end{itemize}
Sufficiency, testing, Bayesian inference, asymptotics, lower bounds, kernels,
and latent models supply the broader statistical foundation. Their full
development need not block independent sections on the direct transformer route.
Advanced expressivity (33) is not required merely to define the model.

\section*{Capstone completion criteria}
\begin{enumerate}[leftmargin=1.8em]
\item Specify a decoder with finite vocabulary, explicit context and termination
conventions, and every map's domain, codomain, and dimensions.
\item Establish a normalized sequence law. For fixed finite length $T$, derive
\[
p_\theta(x_{1:T})=\prod_{t=1}^T p_\theta(x_t\mid x_{<t}).
\]
Variable-length laws require the separately stated termination convention.
\item Derive negative log likelihood and the applicable cross-entropy--KL
decomposition, specifying the data law and weighting.
\item Prove prefix causality, gradients under the chosen regularity conditions,
and cache equivalence in deterministic exact arithmetic.
\item Prove the chosen sampler correct; state arithmetic and memory bounds
in an explicit cost model.
\item Separate established guarantees from empirical scaling, benchmark
performance, and unresolved capability explanations.
\end{enumerate}

\chapter*{Extensions beyond the first milestone}
\addcontentsline{toc}{chapter}{Extensions beyond the first milestone}

These later branches extend the scope; they are not opening prerequisites
or a universal ranking of AI methods.

\section*{A. Learning and decision-making agents}
\begin{enumerate}[leftmargin=1.8em]
\item Markov decision processes and dynamic programming.
\item Partial observation, belief states, and planning.
\item Bandits, exploration, regret, and reinforcement learning.
\item Learned world models, model error, and control.
\item Language-model policies, tools, memory, and interactive tasks.
\end{enumerate}
\textit{Later destination:} a specified agent that learns and acts under
uncertainty, with guarantees in restricted environments. World-model
reinforcement learning gives a concrete research connection~\cite{dreamer}.

\section*{B. Causality, adaptation, and general learning}
Structural causal models; interventions and identification; experimental
design; distribution shift; online and continual learning; multi-agent learning.

\section*{C. Other generative models and modalities}
Graphical models and inference; variational autoencoders; normalizing flows;
score and diffusion models; multimodal learning; graph and geometric models.

\section*{D. Computation and formal reasoning}
Automata, computability, and complexity; formal languages and program semantics;
search and planning; program synthesis; proof-guided learning. These support
longer-term goals involving reasoning and algorithm discovery.

\chapter*{Source checks and open choices}
\addcontentsline{toc}{chapter}{Source checks and open choices}

\textbf{Source use.} This is an original proposed outline. Statistical,
learning-theoretic, and neural coverage was checked against primary book
contents~\cite{wasserman,uml,dl}. Architecture and later topics were checked
against original transformer, language-model, positional, adaptation, and
decoding papers~\cite{attention,gpt,rope,lora,dpo,speculative,yun}.
Checks used contents pages, abstracts, and bibliographic records; they do not
verify the papers' proofs or experiments. Consult precise statements and proof
sources during unit planning. The author's book collection has not yet been
mapped to the plan.

\textbf{Open choices.} Exact probability reference; additional linear-algebra
and analysis boundary; priorities among breadth chapters; routine inference
left implicit in proofs; eventual physical volumes. Revise the plan before
drafting the mathematical definitions that depend on those choices.

\begin{thebibliography}{99}
\raggedright
\bibitem{wasserman} L. Wasserman. \textit{All of Statistics}.
\href{https://link.springer.com/book/10.1007/978-0-387-21736-9}{Publisher's contents}.
\bibitem{uml} S. Shalev-Shwartz and S. Ben-David.
\textit{Understanding Machine Learning: From Theory to Algorithms} (2014).
\href{https://www.cs.huji.ac.il/~shais/UnderstandingMachineLearning/understanding-machine-learning-theory-algorithms.pdf}{Author-hosted text}.
\bibitem{dl} I. Goodfellow, Y. Bengio, and A. Courville.
\textit{Deep Learning} (2016).
\href{https://www.deeplearningbook.org/contents/TOC.html}{Authors' contents}.
\bibitem{attention} A. Vaswani et al. \textit{Attention Is All You Need} (2017).
\href{https://arxiv.org/abs/1706.03762}{arXiv:1706.03762}.
\bibitem{gpt} T. B. Brown et al. \textit{Language Models are Few-Shot Learners} (2020).
\href{https://arxiv.org/abs/2005.14165}{arXiv:2005.14165}.
\bibitem{rope} J. Su et al. \textit{RoFormer: Enhanced Transformer with Rotary Position Embedding} (2021).
\href{https://arxiv.org/abs/2104.09864}{arXiv:2104.09864}.
\bibitem{lora} E. J. Hu et al. \textit{LoRA: Low-Rank Adaptation of Large Language Models} (2021).
\href{https://arxiv.org/abs/2106.09685}{arXiv:2106.09685}.
\bibitem{dpo} R. Rafailov et al. \textit{Direct Preference Optimization:
Your Language Model is Secretly a Reward Model} (2023).
\href{https://arxiv.org/abs/2305.18290}{arXiv:2305.18290}.
\bibitem{speculative} Y. Leviathan, M. Kalman, and Y. Matias.
\textit{Fast Inference from Transformers via Speculative Decoding} (2023).
\href{https://proceedings.mlr.press/v202/leviathan23a.html}{PMLR 202}.
\bibitem{yun} C. Yun et al. \textit{Are Transformers Universal Approximators
of Sequence-to-Sequence Functions?} (2020).
\href{https://arxiv.org/abs/1912.10077}{arXiv:1912.10077}.
\bibitem{dreamer} D. Hafner et al. \textit{Mastering Diverse Domains through
World Models} (2023).
\href{https://arxiv.org/abs/2301.04104}{arXiv:2301.04104}.
\end{thebibliography}
\end{document}
